\section{开发集表达选择与补充机制结果}
\label{sec:supp-development-mechanism}

\subsection{开发集选择的实际工作点}
\label{sec:supp-dev-selected}

开发集选择在五个模型上使用每类4题、共404题，并为IP-VCD、同措辞引导VCD、CDA视觉迁移和同措辞引导M3ID分别比较四个表达。选择目标为$J=(C+TP)/N$；并列时优先整体准确率，再按UNKNOWN、UNCLEAR、UNSURE、I cannot identify it的顺序选择。表\ref{tab:supp-dev-selection}给出实际选择及其2,424题评估结果。同预算的引导VCD和引导M3ID用于分析表达选择的作用。

开发集选择的IP-VCD在Qwen2.5-VL、Qwen3.5-4B、MiniCPM-V2.6和LLaVA-Mistral上，相对原生VCD的评估$J$分别增加15.68、0.91、3.96和14.23个百分点；Gemma-3下降0.66个百分点。Qwen2.5-VL选择UNSURE后$J=44.02\%$，对应评估集四表达最高观测值为$46.37\%$；其余四模型选择结果的$J$与各自最高观测值相同。这些结果给出了仅使用开发输入确定表达后的实际工作点。

\input{statistics/supplementary/dev_selection_table.tex}
\input{statistics/supplementary/dev_eval_counts_table.tex}

表中U、C、S、I依次表示UNKNOWN、UNCLEAR、UNSURE和I cannot identify it。计数C、W、A分别表示正确回答、错误回答和语义弃权；TP与FP分别为参考支持和参考不支持的弃权。每行满足$C+W+A=2,424$且$TP+FP=A$。

\subsection{DoLa与SID算子设置的同题比较}
\label{sec:supp-author-operators}

每类一张图像组成101题配对集合。DoLa采用作者代码的KL方向选择提前层；SID采用作者贪心分支的未截断对比分数和$\alpha=0.5$\cite{chuang2024dola,huo2025sid}。九个DoLa条件与六个SID条件共包含1,515条回复。表\ref{tab:supp-operator101}比较相同输入上的原计算与作者算子计算，表\ref{tab:supp-same101-methods}给出同一集合上各方法的联合效用。此处统计分母为101；完整评估的统计分母为2,424。

DoLa的总正确数从307/909变为306/909，平均$J$变化为$-0.11$个百分点；SID从203/606变为195/606，平均变化为$-1.32$个百分点。模型间变化具有异质性：Gemma-3的DoLa正确数从48增至53，Phi-3.5-Vision从36减至30。Gemma-3在该集合上的DoLa为53/101，IP-VCD为49/101，IP-M3ID为51/101；其余八模型的IP-VCD均高于DoLa。

\input{statistics/supplementary/operator_table.tex}
\input{statistics/supplementary/same101_comparison_table.tex}

\subsection{CDA视觉迁移的条件权重}
\label{sec:supp-cda-weights}

九个模型各选取12条回复路径，覆盖可用的正确、错误、弃权行为与参考阳性、阴性组合，共108条路径。沿实际回复前缀进行五路教师强制测量，得到803个词元位置。采用CDA的无动量校准式\cite{kim2025cda}，令$H_p,H_c$为先验与视觉上下文分支熵，$H_{0p},H_{0c}$为对应空输入熵，则
\begin{align}
 r_p&=\frac{\max(H_p-H_{0p},0)}{H_{0p}},&
 r_c&=\frac{\max(H_c-H_{0c},0)}{H_{0c}},\\
 w_p&=\frac{r_p^2}{r_p+r_c},&
 w_c&=\frac{r_c^2}{r_p+r_c},\qquad w_a=1-w_p-w_c.
\end{align}
当$r_p+r_c=0$时，连续延拓给出$(w_p,w_c,w_a)=(0,0,1)$。三个权重作用于先验、视觉上下文和弃权条件的logit。

803个位置中，418个使用零和连续延拓，145个具有负弃权权重。表\ref{tab:supp-cda-weights}按回复行为给出逐位置均值。Qwen3-VL的正确、错误和弃权三组均出现负的平均$w_a$；MiniCPM-V2.6和OneVision的所测位置均未出现负$w_a$。Gemma-3错误回答组的平均权重幅度达到$10^5$量级，显示熵比值校准可产生强烈的分支放大。CDA在这些路径上表现为随输入和前缀改变的logit线性组合，权重符号与回答行为的对应关系依赖模型。

\input{statistics/supplementary/cda_behavior_table.tex}

这些均值的统计单位为词元位置，路径按行为分层选取。五路分布计算的权重与公式重算值一致；沿相同前缀的当前argmax在776/803个位置与原回复词元相同，其中Gemma-3有26处差异，Qwen3.5-4B有1处差异。表中权重描述本次前缀测量。

\subsection{共享前缀重放中的噪声与硬件条件}
\label{sec:supp-replay-conditions}

六条回复路径上的九个位置在两类硬件和两种会话顺序下共测量36次。正常视觉输入、提示词元、图像和随机种子保持相同；两类硬件生成的扰动视觉张量不同。3090执行在18/18个配对位置上与原回复词元相同，A100或K100执行为0/18。改变会话顺序后，18/18对argmax保持一致（表\ref{tab:supp-replay}）。

进一步在MiniCPM-V2.6的一条苹果派路径上直接共享同一已保存噪声张量，比较两个词元位置和两种会话顺序。3090的4/4个位置保持原词元，A100的首词元在2/2次执行中恢复一致；第二词元在两种顺序下均由句点变为``with''。共享噪声后，该实例的差异集中在第二词元，表明后续解码的数值路径也参与跨硬件差异。

\input{statistics/supplementary/replay_table.tex}
